\emph{Context layers}, curated documentation that an analytics agent
fetches at query time, produce large accuracy gains on text-to-SQL
benchmarks, but a with/without comparison cannot say which part of the
layer does the work. Using a data contract as the instrument, we separate
the knowledge a layer carries from the scaffolding and format that deliver
it, on DABStep and LiveSQLBench, across seven models, with 53 predictions
fixed before their runs. The knowledge does the work: emptying the
contract's prose while keeping its tools and rules costs 26 to 50
hard-task points on every model, and the tools alone are worth $-9$ to
$+15$. What looked like an advantage of delivery was mostly knowledge: two
facts missing from DABStep's manual, stated in words, recover most of the
contract's lead over the pasted manual, and with the facts stated alike
the contract's form adds 4 to 6 points. A contract is worth what its
author could determine: compiled from documentation it is worth the
documentation on both benchmarks; on LiveSQLBench knowledge helps nearly
three times as much where the documentation determines the gold; an author
who can query the data recovers what the data settles and misses the
conventions it cannot reveal. Delivery buys other things. Grown from one
database to 22, the contract in the prompt loses no accuracy up to four
times its size, but looking it up as needed is cheaper from twice its size
(1.7 to 2.4 times at four times) and runs past the context window; the
lookup log traces a wrong answer to the definition behind it; and no
governed run corrupted its database copy, against 50 ungoverned ones. This
version revises our preprint's account, which credited the contract's lead
to delivery.
